\section*{Abstract}
\addcontentsline{toc}{section}{Abstract}

This thesis studies Brzozowski derivatives of regular expressions in two settings,
the classical crisp one and a fuzzy one based on similarity, and verifies the theorems
in the Rocq proof assistant. In the crisp setting it defines derivatives of regular languages and expressions, the
correspondence between the syntactic and the semantic derivative, the finiteness of
the set of derivatives, and the construction of the deterministic automaton that has
canonical derivatives as its states. The original contribution is the fuzzy layer,
which replaces exact symbol matching by a similarity relation cut at a threshold
$\mu$. Its main observation is that, under the minimum triangular norm,
cut-similarity is an equivalence relation. A fuzzy derivative over the alphabet is
then just an ordinary derivative over the quotient alphabet obtained by identifying
similar symbols. Through this quotient construction regularity, the
semantic-syntactic correspondence, finiteness, and correctness of the fuzzy
automaton all carry over from the crisp theory. The automaton recognises exactly
the words that lie within threshold $\mu$ of some word of the original language.
The claims are checked by a proof assistant rather than argued only on paper.

\bigskip
\noindent\textbf{Keywords:} regular expressions, Brzozowski derivatives, finite
automata, fuzzy similarity, triangular norms, formal verification, Rocq.

\clearpage
